\section{Dodgson identities}\label{polyn}

In order to obtain exact results, we need the Dogdson identities in a form which allows to track signs.
They are indeed essential to obtain exact results. In other contexts, the signs can be irrelevant
since the Dodgson polynomials will appear squared or as part of expressions the sign of which is not important,
but it is not the case here. Nevertheless, we want to keep things simple and with
signs which depend on a minimum number of choices in the presentation of the diagram. In particular, the usual
presentation as the determinant of a submatrix with rows and columns removed keeps a dependence on the order
of the rows. The convention we proposed in Section~\ref{presentation} for $U^{I ,J}$ is the determinant of the matrix
$M_G$ with the columns indexed by the set $I$ replaced with ones with only a 1 in the position indexed by the corresponding
element of~$J$. Independence on the order of the edges or the vertices is clear in this case and expansion according to
the columns indexed by~$I$ trivially gives back the usual definition up to a sign, while multilinearity can be used to
show that the same result can be obtained by changing instead the rows indexed by~$J$. Finally, in the case where $I$
and~$J$ have only one element, $U^{i,j}$ only depends on the orientations of $i$ and~$j$.

The sign introduced can be deduced as follows: the elements of $I$ have first to be brought to the beginning of the matrix,
requiring $i_1 - 1$ transpositions to send the first element of $I$ to~1,~${i_2 -2}$ for the following up to $i_k -k $
for the last one if we suppose that the elements of~$I$ are numbered according to the ordering of edges and not the order in~$I$.
We have similar number of transpositions for bringing the elements of $J$ to the front, and the total sign only depends on the parity
of the sum $\sum_{i\in I} i+ \sum_{j\in J} j$. The upper left block of the resulting matrix is a~permutation matrix, which gives a sign
corresponding to the product of the signs associated to the permutations bringing $I$ and~$J$ in the natural order. If one recalls that
the convention in this work is to put a minus sign before one of the incidence block rather than to use the symmetric form used by Schnetz,
one recognizes that there are all the elements of the sign proposed in \cite[equation~(8)]{SC2021}.

Our aim is to express some determinants involving correlators~$U^{i,j}$. Examining the proof of Dodgson
identities, such
as the one found in~\cite{Br2009}, shows that our definition of~$U^{I ,J}$ naturally appears in them. Indeed, $U^{i,j}$
is the $j,i$ component of the adjugate matrix of $M_G$, the matrix which, multiplied by $M_G$ gives $\det(M_G)$
times the identity. We want to express the determinant of the matrix $M$ with elements $U^{i,j}$ such that $i$ is in $I$ and $j$ is in
$J$. We form a matrix of the dimension of~$M_G$ by replacing in the identity matrix the columns with index in~$J$ by the
ones indexed by~$I$ of the adjugate matrix. An expansion along the columns coming from the identity matrix shows
that this new matrix has the same determinant as~$M$.

Multiplying by the left by~$M_G$, one obtains a matrix which has the columns of~$M_G$ except for the columns indexed
by~$J$ which have $\det(M_G)$ in the position indexed by~$I$. We recognize the matrix used to define $U^{J,I}$,
apart from the factors $\det(M_G)$, so that we obtain the required Dodgson identities without any sign ambiguity by
taking the determinants
\[
	\det(M_G) \det\bigl( (U^{i,j})_{i\in I,j\in J} \bigr) = \det(M_G)^{\# J} U^{I ,J}.
\]

